\section{Exercise 3: AM Communication via USRP}
\label{sec:lab2_ex3}

\subsection{Task 1: Functional Understanding of TX and RX Modules}

The final part of the lab studies how the provided USRP transmitter and receiver
modules implement AM communication in practice.

\subsection{USRP Transmitter}

Functionally, the transmitter VI performs the following steps:
\begin{enumerate}
\item open a TX session to the selected USRP device,
\item configure radio parameters such as IQ rate, carrier frequency, gain, and
active antenna,
\item generate the baseband message waveform inside a loop from the selected
message frequency, sample rate, and sample count,
\item apply the modulation index to produce the AM baseband signal,
\item stream the signal continuously using the TX write block so the USRP can
upconvert and transmit it at RF,
\item stop cleanly and close the session when requested.
\end{enumerate}

\begin{figure}[H]
\centering
\labincludegraphics[width=\linewidth]{figures/usrp_tx_vi.png}
\caption{USRP AM transmitter VI block diagram.}
\label{fig:lab2_usrp_tx_vi}
\end{figure}

In this implementation the VI generates the baseband term
\begin{equation}
g(t)=1+\mu\,m(t),
\end{equation}
while the actual RF modulation onto the carrier is performed by the USRP
hardware.

\begin{tcolorbox}[breakable, colback=green!5!white,
colframe=green!60!black,
title={\textbf{Intuition: Baseband Signal Generation in the USRP Transmitter}}]
The general AM signal is
\begin{equation}
s(t)=\left[A_c + A_m m(t)\right]\cos(2\pi f_c t)
=A_c\left[1+\mu\,m(t)\right]\cos(2\pi f_c t).
\end{equation}
In the USRP implementation used here, two simplifications apply:
\begin{itemize}
\item the carrier amplitude is normalized to $A_c=1$ in the baseband signal,
\item the multiplication by $\cos(2\pi f_c t)$ is performed internally by the
USRP hardware rather than explicitly in LabVIEW.
\end{itemize}
Therefore the VI only needs to generate
\begin{equation}
g(t)=1+\mu\,m(t),
\end{equation}
which becomes the real part of the complex IQ stream delivered to the USRP. The
hardware then upconverts that baseband waveform to the configured RF carrier.
\end{tcolorbox}

\subsection{USRP Receiver}

The receiver VI performs:
\begin{enumerate}
\item opening an RX session and configuring the IQ rate, carrier frequency, gain,
and active antenna,
\item starting reception and repeatedly fetching blocks of received samples,
\item applying baseband post-processing such as low-pass filtering, DC removal,
and PSD estimation,
\item displaying the demodulated signal in both time and frequency domains,
\item aborting acquisition and closing the session cleanly on stop.
\end{enumerate}

\begin{figure}[H]
\centering
\labincludegraphics[width=\linewidth]{figures/usrp_rx_vi.png}
\caption{USRP AM receiver VI block diagram.}
\label{fig:lab2_usrp_rx_vi}
\end{figure}

\subsection{How TX and RX Work Together}

The transmitter supplies the USRP with IQ data representing the AM baseband
signal. The hardware then performs RF transmission around the configured carrier
frequency. At the receiving side, the USRP returns baseband samples centered on
that carrier, and the receiver VI completes the recovery through filtering,
offset removal, and visualization.

This is why the modules are mainly configuration, streaming, and baseband
post-processing pipelines rather than full RF implementations in LabVIEW code.

\subsection{Task 2: Single-Tone Transmission and Recovery}

For the single-tone transmission task, the transmitter sends an AM tone and the
receiver output is observed in time and frequency domains.

The expected and reported behaviour is:
\begin{itemize}
\item the demodulated time-domain waveform is sinusoidal, indicating successful
recovery of the message,
\item the demodulated PSD shows a dominant peak at the message frequency, around
$1~\text{kHz}$ in the reference capture.
\end{itemize}

\begin{figure}[H]
\centering
\begin{subfigure}{0.49\linewidth}
\centering
\labincludegraphics[width=\linewidth]{figures/ex3_task2_tx_panel.png}
\caption{USRP TX panel during single-tone transmission.}
\label{fig:lab2_ex3_task2_tx}
\end{subfigure}
\hfill
\begin{subfigure}{0.49\linewidth}
\centering
\labincludegraphics[width=\linewidth]{figures/ex3_task2_rx_panel.png}
\caption{USRP RX panel showing recovered signal and PSD.}
\label{fig:lab2_ex3_task2_rx}
\end{subfigure}
\caption{Exercise 3 Task 2 capture for single-tone transmission and reception.}
\label{fig:lab2_ex3_task2}
\end{figure}

This confirms that the full TX/RX chain is operating correctly: the transmitter
sends an AM-modulated tone and the receiver recovers the baseband message.

\subsection{Task 3: Noise Effect Versus Modulation Index}

To study noise sensitivity, the receiver gain is increased and the modulation
index is varied. The reference observations can be summarized as follows:
\begin{itemize}
\item at $\mu = 1$, the recovered sinusoid is strongest relative to noise and
appears the cleanest,
\item at $\mu = 0.5$, the desired signal amplitude is lower, so noise becomes
more visible,
\item at $\mu = 0.2$, the message component is weakest and waveform perturbations
from noise are clearly noticeable.
\end{itemize}

\begin{figure}[H]
\centering
\begin{subfigure}{0.49\linewidth}
\centering
\labincludegraphics[width=\linewidth]{figures/ex3_task3_mu_1.png}
\caption{$\mu = 1$.}
\label{fig:lab2_ex3_mu1}
\end{subfigure}
\hfill
\begin{subfigure}{0.49\linewidth}
\centering
\labincludegraphics[width=\linewidth]{figures/ex3_task3_mu_05.png}
\caption{$\mu = 0.5$.}
\label{fig:lab2_ex3_mu05}
\end{subfigure}

\vspace{0.5em}

\begin{subfigure}{0.65\linewidth}
\centering
\labincludegraphics[width=\linewidth]{figures/ex3_task3_mu_02.png}
\caption{$\mu = 0.2$.}
\label{fig:lab2_ex3_mu02}
\end{subfigure}
\caption{Exercise 3 Task 3: demodulated output under increased receiver gain for different modulation indices.}
\label{fig:lab2_ex3_task3}
\end{figure}

The reason is that the recovered message amplitude scales with $\mu$, whereas the
receiver noise does not decrease proportionally. As $\mu$ becomes smaller, the
output signal-to-noise ratio deteriorates.

This trend is consistent with the first-order derivation in
Section~\ref{sec:lab2_noise_performance}: for fixed carrier amplitude and
approximately fixed receiver noise level, the output SNR scales roughly as
$\mu^2$. In practical terms, dropping from $\mu=1$ to $\mu=0.5$ corresponds to
an SNR reduction of about $6$ dB, while dropping from $\mu=1$ to $\mu=0.2$
costs about $14$ dB. That scaling explains why the $\mu=0.2$ capture appears
substantially noisier than the $\mu=1$ case.

\subsection{Discussion}

The USRP experiment confirms the same trend already seen in simulation:
\begin{itemize}
\item stronger modulation improves message visibility,
\item weaker modulation makes the recovered baseband increasingly sensitive to
noise,
\item hardware implementation follows the same AM theory, but now the effects are
observed in a real transmission chain rather than only in simulation.
\end{itemize}
